%% LyX 2.4.4 created this file.  For more info, see https://www.lyx.org/.
%% Do not edit unless you really know what you are doing.
\documentclass[english]{article}
\usepackage[T1]{fontenc}
\usepackage[latin9]{inputenc}
\usepackage{geometry}
\geometry{verbose,tmargin=1in,bmargin=1in,lmargin=1in,rmargin=1in}

\makeatletter

%%%%%%%%%%%%%%%%%%%%%%%%%%%%%% LyX specific LaTeX commands.
%% Because html converters don't know tabularnewline
\providecommand{\tabularnewline}{\\}

%%%%%%%%%%%%%%%%%%%%%%%%%%%%%% User specified LaTeX commands.
\usepackage{hyperref}
\usepackage[usenames,dvipsnames]{xcolor} %adds additional color options
\hypersetup{colorlinks=true, urlcolor=Blue}

\usepackage{sectsty} %%one way to change section header font sizes.
\sectionfont{\large}

\usepackage{titlesec}
\titlespacing*{\section}{0pt}{10pt}{0pt}

\usepackage{multicol}
% Added by lyx2lyx
\setlength{\parskip}{\medskipamount}
\setlength{\parindent}{0pt}

\makeatother

\usepackage{babel}
\begin{document}
\begin{center}
{\Large\textbf{Johnson C. Smith University}}\textbf{}\\
\textbf{College Algebra~~|~~MTH 131~~|~~Sections B and I~~|~~Fall
2012}\\
\textbf{\rule[0.25ex]{1\columnwidth}{1pt}}
\par\end{center}

\vspace{-15pt}

\begin{center}
\emph{\textquotedblleft Every act of conscious learning requires the
willingness to suffer an injury to one's self-esteem. That is why
young children, before they are aware of their own self-importance,
learn so easily; and why older persons, especially if vain or important,
cannot learn at all.\textquotedblright{}} \\
\textendash Thomas Szasz, author, professor of psychiatry (1920- )\vspace{-.25cm}
\par\end{center}

\section*{Instructor Information}

\textbf{Name: }Dr.\,James Ashe\\
\textbf{Office location: }307 Davis Hall (back); office hours will
also be held in the Math Lab, EDU 115\\
\textbf{Email: }\href{mailto:jashe@jcsu.edu}{jashe@jcsu.edu}\\
\textbf{Office phone: }(704)\,378-1037\\
\textbf{Office hours: }MWF 10\textendash 11am; MW 1\textendash 4pm;
Tu 12pm\textendash 2pm; F 1pm\textendash 1:30pm, by appointment, or
whenever I\textquoteright m in my office.

\section*{Course Description}

This is a 3 credit hou\textbf{r}\textbf{\emph{ }}course designed to
provide an investigative approach to college algebra using the computer
and graphing calculator as tools. The topics covered include basic
concepts of algebra, functions and graphs, polynomials and rational
functions, exponential and logarithmic functions, systems of equations
and solving systems of equations using matrices, sequences and series.
At least one hour per week of computer assisted laboratory instruction
is required. There are no prerequisites for this course. 

Instructional methods will primarily consist of lectures and in-class
discussions, however other methods may be adopted according to the
particular needs of the class.

\section*{Course Materials }

\textbf{Textbook/MyMathLab: }A textbook/Ebook and a MyMathLab access
code are required. New textbooks and Ebooks usually come packaged
with a MyMathLab code. \\
Textbook: \emph{\href{http://www.amazon.com/College-Algebra-Edition-Robert-Blitzer}{College Algebra, $5^{\mbox{th}}$Edition, Robert Blitzer. ISBN: 0-1314-9264-0}}\\
Alternative options are available here:\\
 \href{http://www.mypearsonstore.com/bookstore/product.asp?isbn=0321559835}{http://www.mypearsonstore.com/bookstore/product.asp?isbn=0321559835}

\textbf{Calculator: }TI-83, TI-83+, or TI-84+\\
Students should bring their calculator and textbook to class.

\section*{Grades}

Grades will be taken from homework, quizzes, and Exams as follows:\textbf{}\\
\textbf{Evaluation: }Quizzes$=20\%$; Exams 1, 2, 3, 4, and the Midterm
Exam $=12\%$ each; Final Exam $=20\%$.

\textbf{Grading scale: $\mbox{A}\geq90>\mbox{B}\geq80>\mbox{C}\geq70>\mbox{D}\geq60>\mbox{F}$}

\textbf{Quizzes: }Quizzes will consist of any graded assignment which
is not an exam. This will include traditional quizzes, MyMathLab assignments,
take home assignments, and in-class projects. The lowest quiz grade
will be dropped. \textbf{}\\
\textbf{Final Exam and Midterm: }Both the final and midterm exam will
be common comprehensive exams given in-class; see the exam schedule
\href{http://www.jcsu.edu/happenings/exam-schedule}{jcsu.edu/happenings/exam-schedule}.
If beneficial, the final exam score will replace the lowest exam score.

\section*{Class Policies}

JCSU policies regarding attendance, academic integrity, grades, and
change of enrollment will be enforced. It is the student\textquoteright s
responsibility to be aware of these policies and of related deadlines
as well as any announcements or syllabus adjustments (written or oral)
made in class. Please see the student handbook for details, \href{http://www.jcsu.edu/docs/handbook.pdf}{jcsu.edu/docs/handbook.pdf}.

\textbf{Attendance: }Class attendance is required for all JCSU students.
Each student is allowed as many hours of absence per term as credit
hours(s) received (not to exceed 4) for the class. Students who exceed
the maximum number of absences may receive a failing grade for the
course. 

Students who may miss classes while representing the University in
an official capacity are exempt from regulations governing absences.
However, absence from class for official University business does
not relieve the student from responsibility for any class assignments
that may be missed during the period of absence.

\textbf{Missed quizzes: }There will be no make-ups for in-class quizzes
unless the class was missed for a University sponsored event, and
advance notice of the absence was given by the student. In this case
the assignment will be excused or a make-up will be given depending
on circumstances. Late take-home quizzes will have 25\% deducted for
each day it is late. \\
\textbf{Missed exams: }If a student misses an exam for a verifiable
excuse, they should send me an e-mail \emph{as soon as possible }so
that a make-up can be arranged. Make-up exams may be significantly
more difficult. If you will be gone for any reason, it will be \emph{your
}responsibility to inform me at least a week in advance to schedule
to take the exam early. 

\textbf{Student support services: }If you need course adaptations
or accommodations for a documented disability please contact the Office
of Disability Services. \\
\href{http://www.jcsu.edu/directory/federal-requirements/general-institutional-information/jcsu-disability-services}{jcsu.edu/directory/federal-requirements/general-institutional-information/jcsu-disability-services},
\\
phone: (704)\,398-1116.

The Office of Counseling Services provides confidential intakes and
assessments, personal counseling and referrals to appropriate resources
to all enrolled students on a walk-in or appointment basis. \\
\href{http://www.jcsu.edu/admissions/future_students/meet_your_counselor}{jcsu.edu/admissions/future\_students/meet\_your\_counselor},
phone (704)\,378-1044.

\section*{Course Content}

\begin{multicols}{2}

\textbf{Chapter 1: Equations and Inequalitie}s \\
1.1 Graphs and Graphing Utilities \\
1.2 Linear Equations and Rational Equations\\
1.3 Models \& Applications \\
1.4 Complex Numbers\\
1.5 Quadratic Functions\\
1.6 Other Types of Equations 

\textbf{Chapter 2: Functions and Graphs} \\
2.1 Basics of Functions and Their Graphs\\
2.3 Linear Functions and Slope \\
2.4 More on Slope

%%\vfill
\columnbreak 

\textbf{Chapter 3: Polynomial and Rational Function}s \\
3.1 Quadratic Functions \\
3.2 Polynomial Functions and Their Graphs 

\textbf{Chapter 4: Exponential and Logarithmic Functions} \\
4.1 Exponential Functions \\
4.2 Logarithmic Functions \\
4.3 Properties of Logarithms 

\end{multicols}

\section*{Important Dates}

\noindent\begin{minipage}[t]{1\columnwidth}%
\begin{tabular}{rcl}
September $21^{\mbox{th}}$ &  & Last day to Add/Drop \tabularnewline
October $5^{\mbox{th}}$ &  & Exam 1\tabularnewline
October $26^{\mbox{th}}$ &  & Exam 2\tabularnewline
November $2^{\mbox{nd}}$ &  & Midterm Exam\tabularnewline
November $30^{\mbox{th}}$ &  & Exam 3\tabularnewline
January $4^{\mbox{th}}$ &  & Exam 4\tabularnewline
January$15^{\mbox{th}}$\textendash $19^{\mbox{th}}$ &  & Common Final Exam\tabularnewline
 &  & \tabularnewline
\end{tabular}%
\end{minipage}

The above exam dates are subject to change. JCSU's Fall 2012 academic
calendar can be found here: \href{http://www.jcsu.edu/happenings/academic-calendar}{jcsu.edu/happenings/academic-calendar},
and the exam schedule can be found here: \href{http://www.jcsu.edu/happenings/exam-schedule}{jcsu.edu/happenings/exam-schedule}
\end{document}
